% Generated draft: requires mathematical and source review.
\documentclass[11pt]{article}
\usepackage[margin=25mm]{geometry}
\usepackage{fontspec}
\setmainfont{Latin Modern Roman}
\usepackage{amsmath,amssymb,mathtools,unicode-math}
\setmathfont{Latin Modern Math}
\usepackage{enumitem,needspace,xurl,hyperref}
\setlength{\parindent}{0pt}
\setlength{\parskip}{5pt}
\setlength{\emergencystretch}{3em}
\newcommand{\N}{\mathbb{N}}
\newcommand{\Z}{\mathbb{Z}}
\newcommand{\Q}{\mathbb{Q}}
\newcommand{\R}{\mathbb{R}}
\newcommand{\C}{\mathbb{C}}
\newcommand{\F}{\mathbb{F}}
\newcommand{\A}{\mathbb{A}}
\newcommand{\PP}{\mathbb{P}}
\newcommand{\E}{\mathbb{E}}
\newcommand{\eps}{\varepsilon}
\newcommand{\dd}{\,\mathrm{d}}
\begin{document}
\section*{UnsolvedMath ID: 800703}
\textbf{Source problem code:} catalogue entry
\subsection{Definitions and mathematical statement}
For a field \(K\), let \(\operatorname{Br}(K)\) denote its Brauer group. For a finite-dimensional central simple \(K\)-algebra \(A\), its period \(\operatorname{per}(A)\) is the order of its class in \(\operatorname{Br}(K)\), and its index \(\operatorname{ind}(A)\) is the degree of the central division algebra representing that class. To make explicit a convention that must be checked against the source, set
\[
b(K)=\inf\bigl\{n\in\Z_{\geq 0}:\operatorname{ind}(A)\mid\operatorname{per}(A)^n\text{ for every finite-dimensional central simple }K\text{-algebra }A\bigr\},
\]
with \(b(K)=\infty\) if no such integer exists. The supplied excerpt does not define its notation \(\operatorname{Br.\!dim}(K)\); in particular, it does not specify whether the period--index bound is imposed over \(K\) alone, as in \(b(K)\), or uniformly over its finite field extensions. This convention remains unresolved.

The requested construction is a notion of field dimension, on a class \(\mathcal K\) of fields to be specified, satisfying, for every \(F\in\mathcal K\),
\[
\operatorname{Br.\!dim}(F)\leq \dim F-1,
\qquad
\dim(F(t))=\dim F+1,
\]
where \(t\) is transcendental over \(F\). To interpret both identities within the same domain of definition, \(\mathcal K\) must also contain \(F(t)\) whenever it contains \(F\). The class \(\mathcal K\), the value set of \(\dim\), and the interpretation of the displayed arithmetic at any infinite values are not specified in the supplied record and are part of the unresolved formulation.

For an integer \(i\geq 0\), a field \(F\) is called a \(C_i\) field if, for every integer \(d\geq 1\), every integer \(N>d^i\), and every homogeneous polynomial of degree \(d\) in \(N\) variables over \(F\), the polynomial has a zero in \(F^N\setminus\{0\}\). The source suggests assigning dimension \(i\) to a \(C_i\) field. Since the \(C_i\) conditions are nested, this suggestion does not by itself define a unique dimension; using the least such \(i\), and handling fields with no such \(i\), would require additional specification.

As an obstruction to using cohomological dimension for this purpose, the supplied source attributes to Merkurjev an example of a field \(F\) with cohomological dimension \(2\) and \(\operatorname{Br.\!dim}(F)=\infty\), citing \([\mathrm{M\!:\!simple}]\). This is retained as a source-reported example; the excerpt does not give the full reference or specify the cohomological-dimension convention.
\subsection{Short English statement}
Construct a notion of dimension for a specified class of fields that bounds Brauer dimension by one less than field dimension and increases by one upon adjoining a transcendental element. The source suggests using \(C_i\) conditions and reports a Merkurjev example obstructing ordinary cohomological dimension. The field class, dimension values, Brauer-dimension convention, and precise meaning of the \(C_i\) suggestion require clarification.
\subsection{Sources}
\begin{itemize}
\item Local UnsolvedMath registry, permanent ID \(800703\), exact source problem code \texttt{catalogue entry}, title: \emph{Open Problems on Central Simple Algebras --- Problem 13}. The supplied record attributes the problem to Auel, Brussel, Garibaldi, and Vishne, with proposed year 2010, and labels it open. This status is supplied registry metadata, not an independently verified literature assessment.
\item Auel, Brussel, Garibaldi, and Vishne, \emph{Open Problems on Central Simple Algebras} (2010), Problem 13, \url{https://arxiv.org/abs/1006.3304}. The supplied provenance records access on 2026-10-16 and extraction by \texttt{arxiv\_latex}; these are supplied provenance claims, not a claim of independent source access or verification.
\item The statement attributes the cohomological-dimension example to Merkurjev through the source reference label \([\mathrm{M\!:\!simple}]\). No full bibliographic entry for that reference is supplied.
\end{itemize}
\end{document}
